\documentclass[10pt,a4paper]{amsart}
\usepackage{amsmath,amssymb,amsthm,graphicx}
% For commutative diagrams we use Xy-pic
\usepackage[all]{xy}

% We use 2cell for 2-commutative diagrams.
\xyoption{2cell}
% Theorem environments.
%
\theoremstyle{plain}
\newtheorem{theorem}[subsection]{Theorem}
\newtheorem{proposition}[subsection]{Proposition}
\newtheorem{lemma}[subsection]{Lemma}

\theoremstyle{definition}
\newtheorem{definition}[subsection]{Definition}
\newtheorem{example}[subsection]{Example}
\newtheorem{exercise}[subsection]{Exercise}
\newtheorem{situation}[subsection]{Situation}

\theoremstyle{remark}
\newtheorem{remark}[subsection]{Remark}
\newtheorem{remarks}[subsection]{Remarks}

\numberwithin{equation}{subsection}

% Macros
%
\def\lim{\mathop{\mathrm{lim}}\nolimits}
\def\colim{\mathop{\mathrm{colim}}\nolimits}
\def\Spec{\mathop{\mathrm{Spec}}}
\def\Hom{\mathop{\mathrm{Hom}}\nolimits}
\def\Ext{\mathop{\mathrm{Ext}}\nolimits}
\def\SheafHom{\mathop{\mathcal{H}\!\mathit{om}}\nolimits}
\def\SheafExt{\mathop{\mathcal{E}\!\mathit{xt}}\nolimits}
\def\Sch{\mathit{Sch}}
\def\Mor{\mathop{\mathrm{Mor}}\nolimits}
\def\Ob{\mathop{\mathrm{Ob}}\nolimits}
\def\Sh{\mathop{\mathit{Sh}}\nolimits}
\def\NL{\mathop{N\!L}\nolimits}
\def\CH{\mathop{\mathrm{CH}}\nolimits}
\def\proetale{{pro\text{-}\acute{e}tale}}
\def\etale{{\acute{e}tale}}
\def\QCoh{\mathit{QCoh}}
\def\Ker{\mathop{\mathrm{Ker}}}
\def\Im{\mathop{\mathrm{Im}}}
\def\Coker{\mathop{\mathrm{Coker}}}
\def\Coim{\mathop{\mathrm{Coim}}}

% Boxtimes
%
\DeclareMathSymbol{\boxtimes}{\mathbin}{AMSa}{"02}

%
% Macros for moduli stacks/spaces
%
\def\QCohstack{\mathcal{QC}\!\mathit{oh}}
\def\Cohstack{\mathcal{C}\!\mathit{oh}}
\def\Spacesstack{\mathcal{S}\!\mathit{paces}}
\def\Quotfunctor{\mathrm{Quot}}
\def\Hilbfunctor{\mathrm{Hilb}}
\def\Curvesstack{\mathcal{C}\!\mathit{urves}}
\def\Polarizedstack{\mathcal{P}\!\mathit{olarized}}
\def\Complexesstack{\mathcal{C}\!\mathit{omplexes}}
% \Pic is the operator that assigns to X its picard group, usage \Pic(X)
% \Picardstack_{X/B} denotes the Picard stack of X over B
% \Picardfunctor_{X/B} denotes the Picard functor of X over B
\def\Pic{\mathop{\mathrm{Pic}}\nolimits}
\def\Picardstack{\mathcal{P}\!\mathit{ic}}
\def\Picardfunctor{\mathrm{Pic}}
\def\Deformationcategory{\mathcal{D}\!\mathit{ef}}
\newenvironment{reference}{}{}
\newenvironment{slogan}{}{}
\newenvironment{history}{}{}
\usepackage{cleveref}
\crefname{originallemma}{original lemma}{original lemmas}
\Crefname{originallemma}{Original lemma}{Original lemmas}
\makeatletter
\newlabel{curves-lemma-nodal-family-etale-local-source}{{0C5D}{source}{}{sourcepointer-1}{}}
\newlabel{curves-lemma-nodal-family-etale-local-source@cref}{{[originallemma][0][]0C5D}{[1][0][]source}{}{}{}}
\makeatother
\crefname{originallemma}{original lemma}{original lemmas}
\Crefname{originallemma}{Original lemma}{Original lemmas}
\makeatletter
\newlabel{curves-lemma-nodal-family-postcompose-etale}{{0CD8}{source}{}{sourcepointer-2}{}}
\newlabel{curves-lemma-nodal-family-postcompose-etale@cref}{{[originallemma][0][]0CD8}{[1][0][]source}{}{}{}}
\makeatother
\crefname{originallemma}{original lemma}{original lemmas}
\Crefname{originallemma}{Original lemma}{Original lemmas}
\makeatletter
\newlabel{lemma-etale-local-at-worst-nodal}{{0CDD}{source}{}{sourcepointer-3}{}}
\newlabel{lemma-etale-local-at-worst-nodal@cref}{{[originallemma][0][]0CDD}{[1][0][]source}{}{}{}}
\makeatother
\crefname{originalsection}{original section}{original sections}
\Crefname{originalsection}{Original section}{Original sections}
\makeatletter
\newlabel{example-blowup}{{0CDC}{source}{}{sourcepointer-4}{}}
\newlabel{example-blowup@cref}{{[originalsection][0][]0CDC}{[1][0][]source}{}{}{}}
\makeatother
\crefname{originallemma}{original lemma}{original lemmas}
\Crefname{originallemma}{Original lemma}{Original lemmas}
\makeatletter
\newlabel{divisors-lemma-flat-base-change-blowing-up}{{0805}{source}{}{sourcepointer-5}{}}
\newlabel{divisors-lemma-flat-base-change-blowing-up@cref}{{[originallemma][0][]0805}{[1][0][]source}{}{}{}}
\makeatother
\crefname{originallemma}{original lemma}{original lemmas}
\Crefname{originallemma}{Original lemma}{Original lemmas}
\makeatletter
\newlabel{algebra-lemma-regular-goes-up}{{07NF}{source}{}{sourcepointer-6}{}}
\newlabel{algebra-lemma-regular-goes-up@cref}{{[originallemma][0][]07NF}{[1][0][]source}{}{}{}}
\makeatother
\crefname{originallemma}{original lemma}{original lemmas}
\Crefname{originallemma}{Original lemma}{Original lemmas}
\makeatletter
\newlabel{morphisms-lemma-smooth-at-point}{{01V9}{source}{}{sourcepointer-7}{}}
\newlabel{morphisms-lemma-smooth-at-point@cref}{{[originallemma][0][]01V9}{[1][0][]source}{}{}{}}
\makeatother
\title{A quasi-compact at-worst-nodal scheme of relative dimension one over a discrete valuation ring is resolved by blowing up singular closed points}
\author{The Stacks Project Authors}
\date{stacks-project:models.tex:local-snapshot-20261012; GNU FDL 1.2 or later}
\begin{document}
\maketitle
\noindent\textit{Editorial note.} The counted unit is the Stacks Project Lemma lemma-blowup-at-worst-nodal at source lines 6190-6206 together with its proof at source lines 6216-6309, a straight adjacent pairing in the chapter file. The chapter is the upstream file models.tex; the pool directory carries the longer name stacks-models and is not used as the chapter name anywhere in this card.\par\medskip
\noindent\textit{Editorial note.} The prose paragraph that sits between the statement and the proof, at source lines 6207-6211, is the author's remark that a slightly stronger statement is also true; it is not part of the counted unit and is not retained. The proof environment itself opens at line 6213, after that paragraph; the first revision had started the retained proof span three lines late, which cut off the opening marker and made the assembler refuse the file.\par\medskip
\noindent\textit{Editorial note.} This card has no context span: the statement and proof are self-contained apart from the seven pointers, and the notions they use are recorded in the standard prerequisites.\par\medskip
\noindent\textit{Editorial note.} The retained proof quotes three results of other chapters, in Curves, Divisors, Algebra and Morphisms, and names two results of this chapter, the étale local structure at a worst nodal point and the worked example of blowing up a node. All are declared as non-load-bearing navigation pointers with their stable Stacks tags looked up in the shared tag index; none of their statements is reproduced inside the span.\par\medskip
\noindent\textit{Editorial note.} The proof contains a commutative diagram of blowing ups drawn with the chapter's Xy-pic setup, which the retained preamble provides; the diagram lies inside the counted proof span and is retained verbatim.\par\medskip
\noindent\textit{Editorial note.} No citation occurs in the retained spans, so no bibliography entry is supplied.\par\medskip
\noindent\textit{Allowed standard prerequisites (editorial).}\begin{itemize}
\item At-worst-nodal schemes of relative dimension one over a discrete valuation ring, and blowing up a closed point. A scheme of relative dimension one over a discrete valuation ring is at-worst-nodal when its singularities are ordinary double points with smooth branches, so that in a neighbourhood of a singular point it is described by a product of a smooth curve germ with a standard node over the ring; the singular points of the special fibre are then finite in number and are exactly the singular points of the total space up to a local computation. Blowing up a closed point replaces it by the projectivised normal cone, and the chapter proves that the blowing up of a singular point remains at-worst-nodal of relative dimension one, that this operation is compatible with the étale local description of the node and with base change to a further blowing up. The author's own lemmas and examples of the same chapter, and the results of the Algebra, Curves, Divisors and Morphisms chapters that the proof quotes, are the allowed inputs.. The counted lemma is the resolution step that the chapter uses to pass from an at-worst-nodal model to a regular one in the semistable reduction section.
\end{itemize}\medskip
\ifdefined\hypertarget\hypertarget{sourcepointer-1}{}\fi
\noindent\textit{Source reference.} Original lemma 0C5D: focused-seeds/AGP/sources/stacks-snapshot-20261011/tags.txt, source lines 15258--15258. Allowed input of the Curves chapter quoted by the retained proof for the étale local description of the family of nodes; the pointer is registered so the reference resolves and the cited statement is not reproduced inside the span. The number is the stable Stacks tag of the target, looked up in the shared tag index.\par\medskip
\ifdefined\hypertarget\hypertarget{sourcepointer-2}{}\fi
\noindent\textit{Source reference.} Original lemma 0CD8: focused-seeds/AGP/sources/stacks-snapshot-20261011/tags.txt, source lines 15533--15533. Allowed input of the Curves chapter quoted by the retained proof for the compatibility of that local description with postcomposition by an étale map; the pointer is registered so the reference resolves and the cited statement is not reproduced inside the span. The number is the stable Stacks tag of the target, looked up in the shared tag index.\par\medskip
\ifdefined\hypertarget\hypertarget{sourcepointer-3}{}\fi
\noindent\textit{Source reference.} Original lemma 0CDD: focused-seeds/AGP/sources/stacks-snapshot-20261011/tags.txt, source lines 15538--15538. Navigation pointer to the chapter's own statement of the étale local structure at a worst nodal point, which the retained proof uses; the cited statement is an allowed input of the same section and its proof is not reproduced here. The number is the stable Stacks tag of the target, looked up in the shared tag index.\par\medskip
\ifdefined\hypertarget\hypertarget{sourcepointer-4}{}\fi
\noindent\textit{Source reference.} Original section 0CDC: focused-seeds/AGP/sources/stacks-snapshot-20261011/tags.txt, source lines 15537--15537. Navigation pointer to the chapter's own worked example of the blowing up of a node, used by the retained proof to recognise the local model. The assembler's whitelist of kinds does not contain the kind that corresponds to an example, so following the established mapping this stub is declared with the kind section; the cited example is an allowed input and its text is not reproduced here. The number is the stable Stacks tag 0CDC, looked up in the shared tag index.\par\medskip
\ifdefined\hypertarget\hypertarget{sourcepointer-5}{}\fi
\noindent\textit{Source reference.} Original lemma 0805: focused-seeds/AGP/sources/stacks-snapshot-20261011/tags.txt, source lines 10170--10170. Allowed input of the Divisors chapter quoted by the retained proof for the flatness under base change of a blowing up; the pointer is registered so the reference resolves and the cited statement is not reproduced inside the span. The number is the stable Stacks tag of the target, looked up in the shared tag index.\par\medskip
\ifdefined\hypertarget\hypertarget{sourcepointer-6}{}\fi
\noindent\textit{Source reference.} Original lemma 07NF: focused-seeds/AGP/sources/stacks-snapshot-20261011/tags.txt, source lines 9760--9760. Allowed input of the Algebra chapter quoted by the retained proof for the behaviour of regularity under the local ring extension produced by the blowup; the pointer is registered so the reference resolves and the cited statement is not reproduced inside the span. The number is the stable Stacks tag of the target, looked up in the shared tag index.\par\medskip
\ifdefined\hypertarget\hypertarget{sourcepointer-7}{}\fi
\noindent\textit{Source reference.} Original lemma 01V9: focused-seeds/AGP/sources/stacks-snapshot-20261011/tags.txt, source lines 2438--2438. Allowed input of the Morphisms chapter quoted by the retained proof for the smoothness at a point of the total space away from the singular points; the pointer is registered so the reference resolves and the cited statement is not reproduced inside the span. The number is the stable Stacks tag of the target, looked up in the shared tag index.\par\medskip
\begin{lemma}
\label{lemma-blowup-at-worst-nodal}
Let $R$ be a discrete valuation ring. Let $X$ be a quasi-compact scheme which
is at-worst-nodal of relative dimension $1$ with smooth generic fibre over $R$.
Then there exists $m \geq 0$ and a sequence
$$
X_m \to \ldots \to X_1 \to X_0 = X
$$
such that
\begin{enumerate}
\item $X_{i + 1} \to X_i$ is the blowing up of a closed point
$x_i$ where $X_i$ is singular,
\item $X_i \to \Spec(R)$ is at-worst-nodal of relative dimension $1$,
\item $X_m$ is regular.
\end{enumerate}
\end{lemma}


\begin{proof}
Since $X$ is quasi-compact we see that the special fibre $X_k$ is quasi-compact.
Since the singularities of $X_k$ are at-worst-nodal, we see
that $X_k$ has a finite number of nodes and is otherwise
smooth over $k$. As $X \to \Spec(R)$ is flat with smooth generic
fibre it follows that $X$ is smooth over $R$ except at the
finite number of nodes of $X_k$
(use Morphisms, Lemma \ref{morphisms-lemma-smooth-at-point}).
It follows that $X$ is regular at every point except for possibly
the nodes of its special fibre
(see Algebra, Lemma \ref{algebra-lemma-regular-goes-up}).
Let $x \in X$ be such a node.
Choose a diagram
$$
\xymatrix{
X \ar[d] &
U \ar[r] \ar[d] \ar[l] &
\Spec(A) \ar[dl] \\
\Spec(R) &
\Spec(R') \ar[l]
}
$$
as in Lemma \ref{lemma-etale-local-at-worst-nodal}.
Observe that the case $A = R'[u, v]/(uv)$ cannot
occur, as this would mean that the generic fibre of
$X/R$ is singular (tiny detail omitted). Thus $A = R'[u, v]/(uv - \pi^n)$
for some $n \geq 0$. Since $x$ is a singular point,
we have $n \geq 2$, see discussion in
Example \ref{example-blowup}.

\medskip\noindent
After shrinking $U$ we may assume there is
a unique point $u \in U$ mapping to $x$.
Let $w \in \Spec(A)$ be the image of $u$.
We may also assume that $u$ is the unique point of $U$
mapping to $w$.
Since the two horizontal arrows are \'etale
we see that $u$, viewed as a closed subscheme of $U$,
is the scheme theoretic inverse image of $x \in X$ and the
scheme theoretic inverse image of $w \in \Spec(A)$.
Since blowing up commutes with flat base change
(Divisors, Lemma \ref{divisors-lemma-flat-base-change-blowing-up})
we find a commutative diagram
$$
\xymatrix{
X' \ar[d] &
U' \ar[l] \ar[d] \ar[r] &
W' \ar[d] \\
X & U \ar[l] \ar[r] & \Spec(A)
}
$$
with cartesian squares where the vertical arrows are the blowing
up of $x, u, w$ in $X, U, \Spec(A)$. The scheme $W'$ was described
in Example \ref{example-blowup}. We saw there that $W'$
at-worst-nodal of relative dimension $1$ over $R'$. Thus
$W'$ is at-worst-nodal of relative dimension $1$ over $R$
(Algebraic Curves, Lemma \ref{curves-lemma-nodal-family-postcompose-etale}).
Hence $U'$ is at-worst-nodal of relative dimension $1$ over $R$ (see
Algebraic Curves, Lemma \ref{curves-lemma-nodal-family-etale-local-source}).
Since $X' \to X$ is an isomorphism over the complement of $x$,
we conclude the same thing is true of $X'/R$ (by
Algebraic Curves, Lemma \ref{curves-lemma-nodal-family-etale-local-source}
again).

\medskip\noindent
Finally, we need to argue that after doing a finite number
of these blowups we arrive at a regular model $X_m$.
This is rather clear because the ``invariant'' $n$ decreases by $2$
under the blowup described above, see computation in
Example \ref{example-blowup}.
However, as we want to avoid precisely defining this invariant
and establishing its properties, we in stead argue as follows.
If $n = 2$, then $W'$ is regular and hence $X'$
is regular at all points lying over $x$ and we have
decreased the number of singular points of $X$ by $1$.
If $n > 2$, then the unique singular point $w'$ of $W'$ lying over $w$
has $\kappa(w) = \kappa(w')$. Hence $U'$ has a unique
singular point $u'$ lying over $u$ with $\kappa(u) = \kappa(u')$.
Clearly, this implies that $X'$ has a unique singular point $x'$
lying over $x$, namely the image of $u'$. Thus we can
argue exactly as above that we get a commutative diagram
$$
\xymatrix{
X'' \ar[d] &
U'' \ar[l] \ar[d] \ar[r] &
W'' \ar[d] \\
X' & U' \ar[l] \ar[r] & W'
}
$$
with cartesian squares where the vertical arrows are the blowing
up of $x', u', w'$ in $X', U', W'$. Continuing like this we get
a compatible sequence of blowups which stops after
$\lfloor n/2 \rfloor$ steps. At the completion of this process
the scheme $X^{(\lfloor n/2 \rfloor)}$ will have one fewer
singular point than $X$. Induction on the number of singular
points completes the proof.
\end{proof}

\end{document}
